\documentclass[10pt,a4paper]{article}
\usepackage[margin=0.64in]{geometry}
\usepackage[T1]{fontenc}
\usepackage{lmodern}
\usepackage{microtype}
\usepackage{enumitem}
\usepackage[hidelinks]{hyperref}
\usepackage{xcolor}
\usepackage{titlesec}
\setlength{\parindent}{0pt}
\setlength{\parskip}{2pt}
\setlist[itemize]{leftmargin=1.4em,topsep=1pt,itemsep=0.6pt,parsep=0pt}
\titleformat{\section}{\large\bfseries}{}{0em}{}[\titlerule]
\titlespacing*{\section}{0pt}{0.55em}{0.22em}
\titlespacing*{\subsection}{0pt}{0.28em}{0.1em}
\pagestyle{empty}
\begin{document}
{\LARGE\bfseries Zhiheng Zhang}\hfill Melbourne, Australia\\
\href{mailto:zhiheng0mera@gmail.com}{zhiheng0mera@gmail.com}\quad
\href{mailto:zhihezhang@student.unimelb.edu.au}{zhihezhang@student.unimelb.edu.au}\\
\href{https://github.com/zhiheng-zhang-Mera}{github.com/zhiheng-zhang-Mera}

\section{Research Profile}
Mobile and ubiquitous systems researcher developing reliable, goal-grounded personal AI execution across lightweight wearable interfaces and nearby computing infrastructure.

\section{Education}
\textbf{University of Melbourne}, Australia \hfill 2025--2026 (expected)\\
Master of Computer Science (research project underway)

\textbf{University of British Columbia}, Canada \hfill Completed 2024\\
Bachelor of Science in Computer Science

\section{Selected Research Engineering Projects}
\subsection*{Utopia / Digital City \hfill \small Bounded multi-device execution and evaluation substrate}
\href{https://github.com/zhiheng-zhang-Mera/utopia}{github.com/zhiheng-zhang-Mera/utopia}
\begin{itemize}
\item Built Web and Android control surfaces for a bounded personal-computing system with authenticated device routing and recoverable execution state.
\item Verified a three-end setup with two physical Windows execution workers and Android control; PCF integration includes bounded real cross-host task results and explicit resource-observation contracts.
\item Maintained research-evaluation evidence for fault injection, replay, traceable results and independent bounded reproducibility; wearable deployment and full dual-server failover remain future work.
\end{itemize}

\subsection*{DS-Hns \hfill \small Long-running software-agent runtime}
\href{https://github.com/zhiheng-zhang-Mera/DS-Hns}{github.com/zhiheng-zhang-Mera/DS-Hns}
\begin{itemize}
\item Developed durable task lifecycle and queueing for repository-scale AI-assisted work, including restart/resume recovery and failure-isolated extensions.
\item Used explicit terminal-state and recovery evidence to study the gap between continued process execution and genuine task completion.
\end{itemize}

\subsection*{Codex Boss \hfill \small Evidence-governed multi-agent execution}
\href{https://github.com/zhiheng-zhang-Mera/Codex-Boss}{github.com/zhiheng-zhang-Mera/Codex-Boss}
\begin{itemize}
\item Built orchestration and verification foundations for multi-agent workflows, including acceptance records, critique and bounded completion claims.
\item Preserved unresolved negative cases and traceable failures rather than treating superficial success signals as proof of completion.
\end{itemize}

\section{Research Methods and Skills}
System prototyping; multi-device software integration; long-running task lifecycles; fault injection; reproducible experiment and evidence design; test and failure analysis; independent cross-host review.

\section{Proposed Doctoral Research}
How an agent can preserve authorized intent and recover a personal task when wireless connectivity changes between wearable control surfaces and a stable personal compute hub. Experimental priority: Run real two-host-plus-mobile-client tasks under controlled network loss and latency variation; compare recovery and status correctness with fixed-plan approaches, distinguishing emulated from physical wireless evidence.

\end{document}
